% Draft notes for gradual typing, built as their own document (graduality.tex) rather than as
% a chapter of the specification.
\section{Graduality}
\label{sec:graduality}

\figref{graduality-precision} gives the gradual types. Define $\TypeSet$ to be the set of
types, as given in \specification, and $\GTypeSet$ to be the set
of gradual types; $\sigma$ and $\tau$ range over $\TypeSet$, and $\gty{\sigma}$ and $\gty{\tau}$ over
$\GTypeSet$. Concretisation and abstraction are from \citet{garcia2016}; precision is from
\citet{siek2015}.

\begin{figure}
\small
%
\begin{center}
\begin{tabular}{rcll}
\grammarGroup{Gradual types} \\
$\gty{\tau}$ & $::=$ & $\tyAny$ & (Any) \\
    & $\mid$ & $\nu$ & (Primitive) \\
    & $\mid$ & $\tyList{\gty{\tau}}$ & (List) \\
    & $\mid$ & $\tyTuple{\gty{\tau}_1, \ldots, \gty{\tau}_n}$ & (Tuple) \\
    & $\mid$ & $\tyDict{\gty{\tau}}$ & (Dict) \\
    & $\mid$ & $\tyCallable{\gty{\tau}_1, \ldots, \gty{\tau}_n}{\gty{\tau}}$ & (Callable) \\
    & $\mid$ & $\tyLiteral{\ell}$ & (Literal) \\
    & $\mid$ & $C$ & (Class) \\
    & $\mid$ & $\tyUnion{\gty{\sigma}}{\gty{\tau}}$ & (Union)
\end{tabular}
\end{center}

\vspace{2mm}

\flushleft \shadebox{$\gty{\sigma} \precTy \gty{\tau}$}
\begin{mathpar}
\small
\inferrule*[lab={\ruleName{ref-any}}]
{
  \strut
}
{
  \gty{\tau} \precTy \tyAny
}
\and
\inferrule*[lab={\ruleName{ref-prim}}]
{
  \strut
}
{
  \nu \precTy \nu
}
\and
\inferrule*[lab={\ruleName{ref-literal}}]
{
  \strut
}
{
  \tyLiteral{\ell} \precTy \tyLiteral{\ell}
}
\and
\inferrule*[lab={\ruleName{ref-class}}]
{
  \strut
}
{
  C \precTy C
}
\and
\inferrule*[lab={\ruleName{ref-list}}]
{
  \gty{\sigma} \precTy \gty{\tau}
}
{
  \tyList{\gty{\sigma}} \precTy \tyList{\gty{\tau}}
}
\and
\inferrule*[lab={\ruleName{ref-dict}}]
{
  \gty{\sigma} \precTy \gty{\tau}
}
{
  \tyDict{\gty{\sigma}} \precTy \tyDict{\gty{\tau}}
}
\and
\inferrule*[lab={\ruleName{ref-tuple}}]
{
  \gty{\sigma}_i \precTy \gty{\tau}_i \quad (\forall i)
}
{
  \tyTuple{\gty{\sigma}_1, \ldots, \gty{\sigma}_n} \precTy \tyTuple{\gty{\tau}_1, \ldots, \gty{\tau}_n}
}
\and
\inferrule*[lab={\ruleName{ref-callable}}]
{
  \gty{\sigma}_i \precTy \gty{\tau}_i \quad (\forall i) \\
  \gty{\sigma} \precTy \gty{\tau}
}
{
  \tyCallable{\gty{\sigma}_1, \ldots, \gty{\sigma}_n}{\gty{\sigma}} \precTy \tyCallable{\gty{\tau}_1, \ldots, \gty{\tau}_n}{\gty{\tau}}
}
\and
\inferrule*[lab={\ruleName{ref-union}}]
{
  \gty{\sigma} \precTy \gty{\tau} \\
  \gty{\sigma}' \precTy \gty{\tau}'
}
{
  \tyUnion{\gty{\sigma}}{\gty{\sigma}'} \precTy \tyUnion{\gty{\tau}}{\gty{\tau}'}
}
\end{mathpar}
\caption{Gradual types and refinement}
\label{fig:graduality-precision}
\end{figure}


\begin{definition}[Precision]
\label{def:precision}
Precision $\precTy$, a partial order on gradual types, is defined in \figref{graduality-precision}.
\end{definition}

\begin{lemma}[Discreteness of static types]
\label{lem:static-discrete}
If $\sigma \precTy \tau$ then $\sigma = \tau$.
\end{lemma}

\begin{definition}[Concretisation]
\label{def:concretisation}
Concretisation $\gamma : \GTypeSet \to \powersetplus{\TypeSet}$ is defined in \figref{graduality-concretisation}.
\end{definition}

\begin{figure}
\small
\begin{defalign}
\gamma(\tyAny) &= \TypeSet \\*
\gamma(\nu) &= \set{\nu} \\*
\gamma(\tyList{\gty{\tau}}) &= \set{\tyList{\tau} \mid \tau \in \gamma(\gty{\tau})} \\*
\gamma(\tyTuple{\seq{\gty{\tau}}}) &= \set{\tyTuple{\seq{\tau}} \mid \seq{\tau} \in \gamma(\seq{\gty{\tau}})} \\*
\gamma(\tyDict{\gty{\tau}}) &= \set{\tyDict{\tau} \mid \tau \in \gamma(\gty{\tau})} \\*
\gamma(\tyCallable{\seq{\gty{\tau}}}{\gty{\tau}}) &= \set{\tyCallable{\seq{\tau}}{\tau} \mid \seq{\tau} \in \gamma(\seq{\gty{\tau}}), \; \tau \in \gamma(\gty{\tau})} \\*
\gamma(\tyLiteral{\ell}) &= \set{\tyLiteral{\ell}} \\*
\gamma(C) &= \set{C} \\*
\gamma(\tyUnion{\gty{\sigma}}{\gty{\tau}}) &= \set{(\tyUnion{\sigma}{\tau}) \mid \sigma \in \gamma(\gty{\sigma}), \; \tau \in \gamma(\gty{\tau})} \\*[2mm]
\gamma(\seq{\gty{\tau}}) &= \gamma(\gty{\tau}_1) \times \cdots \times \gamma(\gty{\tau}_n)
\end{defalign}
\caption{Concretisation}
\label{fig:graduality-concretisation}
\end{figure}


\begin{lemma}[Preservation and reflection of precision]
\label{lem:precision-concretisation}
$\gty{\sigma} \precTy \gty{\tau}$ if and only if $\gamma(\gty{\sigma}) \subseteq \gamma(\gty{\tau})$.
\end{lemma}

A gradual type $\gty{\tau}$ \emph{approximates} a set $S$ of types when $S \subseteq \gamma(\gty{\tau})$;
$\tyAny$ approximates every set of types. A least approximant is a best abstraction in the sense
of \citet{cousot1977}. Preservation is monotonicity of concretisation. Reflection forces every
approximant of $S$ to lie above the join of $S$, so the join is the least approximant; without
reflection it would approximate $S$ but need not be least.

\begin{lemma}[Least approximant]
\label{lem:least-approximant}
For every nonempty set $S$ of types, the least upper bound $\bigsqcup S$ under $\precTy$ exists and
is the least $\gty{\tau}$ that approximates $S$.
\end{lemma}

\begin{proof}
Since $\gamma(\sigma) = \set{\sigma}$ for every type $\sigma$, approximating $S$ is the same as
bounding $S$ above under $\precTy$ (Lemma~\ref{lem:precision-concretisation}). The join of two
gradual types (\figref{graduality-join}) is their least upper bound, by induction on the pair, so
every finite nonempty set of gradual types has a least upper bound, by induction on its size. Fix
$\sigma \in S$. The gradual types above $\sigma$ are obtained by replacing subterms of $\sigma$ with
$\tyAny$, so there are finitely many; the joins $\bigsqcup F$ for finite $F \subseteq S$ containing
$\sigma$ lie among them and increase with $F$, so some $\bigsqcup F$ is maximal. Then
$\tau \precTy \bigsqcup (F \cup \set{\tau}) = \bigsqcup F$ for every $\tau \in S$ by maximality, so
$\bigsqcup F$ bounds $S$; and every bound of $S$ bounds $F$, so $\bigsqcup F$ is least.
\end{proof}

\begin{figure}
\small
\begin{defalign}
\tyAny \gjoin \gty{\tau} &= \tyAny \\*
\gty{\tau} \gjoin \tyAny &= \tyAny \\*
\nu \gjoin \nu &= \nu \\*
\tyList{\gty{\sigma}} \gjoin \tyList{\gty{\tau}} &= \tyList{\gty{\sigma} \gjoin \gty{\tau}} \\*
\tyTuple{\seq{\gty{\sigma}}} \gjoin \tyTuple{\seq{\gty{\tau}}} &= \tyTuple{\gty{\sigma}_1 \gjoin \gty{\tau}_1, \ldots, \gty{\sigma}_n \gjoin \gty{\tau}_n} \\*
\tyDict{\gty{\sigma}} \gjoin \tyDict{\gty{\tau}} &= \tyDict{\gty{\sigma} \gjoin \gty{\tau}} \\*
\tyCallable{\seq{\gty{\sigma}}}{\gty{\sigma}} \gjoin \tyCallable{\seq{\gty{\tau}}}{\gty{\tau}} &= \tyCallable{\gty{\sigma}_1 \gjoin \gty{\tau}_1, \ldots, \gty{\sigma}_n \gjoin \gty{\tau}_n}{\gty{\sigma} \gjoin \gty{\tau}} \\*
\tyLiteral{\ell} \gjoin \tyLiteral{\ell} &= \tyLiteral{\ell} \\*
C \gjoin C &= C \\*
(\tyUnion{\gty{\sigma}}{\gty{\sigma}'}) \gjoin (\tyUnion{\gty{\tau}}{\gty{\tau}'}) &= \tyUnion{(\gty{\sigma} \gjoin \gty{\tau})}{(\gty{\sigma}' \gjoin \gty{\tau}')} \\*
\gty{\sigma} \gjoin \gty{\tau} &= \tyAny && \quad \text{otherwise}
\end{defalign}
\caption{Join}
\label{fig:graduality-join}
\end{figure}


\begin{definition}[Abstraction]
\label{def:abstraction}
Abstraction $\alpha : \powersetplus{\TypeSet} \to \GTypeSet$ is defined by
$\alpha(S) = \bigsqcup S$ (Lemma~\ref{lem:least-approximant}).
\end{definition}

\noindent Every gradual type approximates the empty set, and $\GTypeSet$ has no least element under
$\precTy$, so the empty set is excluded from the domain of $\alpha$. Concretisation and abstraction
form a \emph{Galois insertion}: $\alpha(S) \precTy \gty{\tau}$ if and only if
$S \subseteq \gamma(\gty{\tau})$ (Lemma~\ref{lem:least-approximant}), and
$\alpha(\gamma(\gty{\tau})) = \gty{\tau}$ by reflection
(Lemma~\ref{lem:precision-concretisation}).

A relation on types lifts to gradual types through concretisation \citep{garcia2016}: the lifted
relation holds when the concretisations contain a related pair. Subtyping and equivalence on types
are as defined in \specification; equivalence lifts to \emph{consistency}.

\begin{definition}[Consistency]
\label{def:consistency}
$\Sigma \vdash \gty{\sigma} \gtyEquiv \gty{\tau}$ if and only if $\Sigma \vdash \sigma \tyEquiv \tau$ for some
$\sigma \in \gamma(\gty{\sigma})$ and $\tau \in \gamma(\gty{\tau})$.
\end{definition}

\figref{graduality-subtyping} defines \emph{consistent subtyping} by lifting the subtyping rules
of \specification rule by rule: types become gradual types, subtyping premises become consistent
subtyping, equivalence premises become consistency, and two axioms place $\tyAny$ below and above
every gradual type. The intended characterisation is the lifting of subtyping itself; the rules
satisfy one direction.

\begin{figure}
\flushleft \shadebox{$\Sigma \vdash \gty{\sigma} \gsubty \gty{\tau}$}
\begin{mathpar}
\small
\inferrule*[lab={\ruleName{gsubty-any-left}}]
{
  \strut
}
{
  \Sigma \vdash \tyAny \gsubty \gty{\tau}
}
\and
\inferrule*[lab={\ruleName{gsubty-any-right}}]
{
  \strut
}
{
  \Sigma \vdash \gty{\sigma} \gsubty \tyAny
}
\and
\inferrule*[lab={\ruleName{gsubty-refl}}]
{
  \strut
}
{
  \Sigma \vdash \gty{\tau} \gsubty \gty{\tau}
}
\and
\inferrule*[lab={\ruleName{gsubty-class}}]
{
  D \in \ancestors(\Sigma, C)
}
{
  \Sigma \vdash C \gsubty D
}
\and
\inferrule*[lab={\ruleName{gsubty-never}}]
{
  \strut
}
{
  \Sigma \vdash \tyNever \gsubty \gty{\tau}
}
\and
\inferrule*[lab={\ruleName{gsubty-object}}]
{
  \strut
}
{
  \Sigma \vdash \gty{\tau} \gsubty \tyObject
}
\and
\inferrule*[lab={\ruleName{gsubty-int-float}}]
{
  \strut
}
{
  \Sigma \vdash \tyInt \gsubty \tyFloat
}
\and
\inferrule*[lab={\ruleName{gsubty-sized}}]
{
  \gty{\tau} \text{ a list, dictionary, string or tuple type}
}
{
  \Sigma \vdash \gty{\tau} \gsubty \tySized
}
\and
\inferrule*[lab={\ruleName{gsubty-list}}]
{
  \Sigma \vdash \gty{\sigma} \gtyEquiv \gty{\tau}
}
{
  \Sigma \vdash \tyList{\gty{\sigma}} \gsubty \tyList{\gty{\tau}}
}
\and
\inferrule*[lab={\ruleName{gsubty-dict}}]
{
  \Sigma \vdash \gty{\sigma} \gtyEquiv \gty{\tau}
}
{
  \Sigma \vdash \tyDict{\gty{\sigma}} \gsubty \tyDict{\gty{\tau}}
}
\and
\inferrule*[lab={\ruleName{gsubty-callable}}]
{
  \Sigma \vdash \gty{\tau}_i \gsubty \gty{\sigma}_i \quad (\forall i) \\
  \Sigma \vdash \gty{\sigma} \gsubty \gty{\tau}
}
{
  \Sigma \vdash \tyCallable{\gty{\sigma}_1, \ldots, \gty{\sigma}_n}{\gty{\sigma}} \gsubty \tyCallable{\gty{\tau}_1, \ldots, \gty{\tau}_n}{\gty{\tau}}
}
\and
\inferrule*[lab={\ruleName{gsubty-tuple}}]
{
  \Sigma \vdash \gty{\sigma}_i \gsubty \gty{\tau}_i \quad (\forall i)
}
{
  \Sigma \vdash \tyTuple{\gty{\sigma}_1, \ldots, \gty{\sigma}_n} \gsubty \tyTuple{\gty{\tau}_1, \ldots, \gty{\tau}_n}
}
\and
\inferrule*[lab={\ruleName{gsubty-literal}}]
{
  \Sigma \vdash \baseType(\ell) \gsubty \gty{\tau}
}
{
  \Sigma \vdash \tyLiteral{\ell} \gsubty \gty{\tau}
}
\and
\inferrule*[lab={\ruleName{gsubty-union-left}}]
{
  \Sigma \vdash \gty{\sigma} \gsubty \gty{\tau}
}
{
  \Sigma \vdash \gty{\sigma} \gsubty \tyUnion{\gty{\tau}}{\gty{\tau}'}
}
\and
\inferrule*[lab={\ruleName{gsubty-union-right}}]
{
  \Sigma \vdash \gty{\sigma} \gsubty \gty{\tau}'
}
{
  \Sigma \vdash \gty{\sigma} \gsubty \tyUnion{\gty{\tau}}{\gty{\tau}'}
}
\and
\inferrule*[lab={\ruleName{gsubty-union}}]
{
  \Sigma \vdash \threeSub{\gty{\tau}}{\gty{\sigma}}{\gty{\tau}'}
}
{
  \Sigma \vdash \tyUnion{\gty{\tau}}{\gty{\tau}'} \gsubty \gty{\sigma}
}
\end{mathpar}
\caption{Consistent subtyping: lifted rules}
\label{fig:graduality-subtyping}
\end{figure}


\begin{lemma}[Completeness of the lifted rules]
\label{lem:consistent-subtyping-complete}
If $\Sigma \vdash \sigma \subty \tau$ for some $\sigma \in \gamma(\gty{\sigma})$ and
$\tau \in \gamma(\gty{\tau})$ then $\Sigma \vdash \gty{\sigma} \gsubty \gty{\tau}$.
\end{lemma}

The converse fails at \ruleName{gsubty-union} alone. The premises of \ruleName{gsubty-union} share
their right-hand type, for which the lifting of subtyping takes one witness,
whereas the lifted premises take witnesses independently. The rules derive
$\Sigma \vdash \tyUnion{\tyList{\tyInt}}{\tyList{\tyStr}} \gsubty \tyList{\tyAny}$, since
$\tyInt \gtyEquiv \tyAny$ and $\tyStr \gtyEquiv \tyAny$; but
$\gamma(\tyList{\tyAny}) = \set{\tyList{\tau} \mid \tau \in \TypeSet}$, and
$\Sigma \vdash \tyUnion{\tyList{\tyInt}}{\tyList{\tyStr}} \subty \tyList{\tau}$ requires
$\tyInt \tyEquiv \tau$ and $\tyStr \tyEquiv \tau$.

\subsection{Joint lifting of the union premises}

The premises of \ruleName{gsubty-union} share their right-hand type, so they must be lifted
jointly, as one predicate over three gradual types.

\begin{definition}[$\USub$]
\label{def:usub}
$\Sigma \vdash \USub(\gty{\tau}, \gty{\tau}'; \gty{\sigma})$ if and only if
$\Sigma \vdash \tau \subty \sigma$ and $\Sigma \vdash \tau' \subty \sigma$ for some
$\tau \in \gamma(\gty{\tau})$, $\tau' \in \gamma(\gty{\tau}')$ and $\sigma \in \gamma(\gty{\sigma})$.
\end{definition}

\noindent The predicate is symmetric in its first two arguments. Replacing the premises of
\ruleName{gsubty-union} with $\USub$ alone,
%
\begin{mathpar}
\inferrule*[lab={\ruleName{gsubty-union}}]
{
  \Sigma \vdash \USub(\gty{\tau}, \gty{\tau}'; \gty{\sigma})
}
{
  \Sigma \vdash \tyUnion{\gty{\tau}}{\gty{\tau}'} \gsubty \gty{\sigma}
}
\end{mathpar}
%
makes the rule match the lifting of subtyping by construction, because a concrete union lies below
$\sigma$ exactly when both components do; the counterexample above is no longer derivable.

An inductive characterisation of $\USub$ can be extracted from a decidability proof: each case of
the induction either exhibits witnesses or refutes their existence, and the positive cases, read
as inference rules, give \figref{graduality-usub}. The proof is a device for finding the rules;
the surviving obligation, that the rules derive the predicate exactly, follows from the same case
analysis, because every positive case is an equivalence.

\begin{theorem}[Decidability of $\USub$]
\label{thm:usub-decidable}
The predicate of Definition~\ref{def:usub} is decidable.
\end{theorem}

\begin{proof}[Proof sketch]
By induction on the triple, ordered by the total size of the three types and then by the size of
the first; cases on the triple. A left type $\tyAny$ or $\tyNever$ takes the witness $\tyNever$,
leaving a consistent subtyping check on the other left type (\ruleName{usub-any-left},
\ruleName{usub-never-left}). A left union rebrackets, because its witness is a union, below
$\sigma$ exactly when both components are (\ruleName{usub-union}). Otherwise both left witnesses
are constructor-headed. On the right, $\tyAny$ admits the witness $\tyObject$
(\ruleName{usub-any-right}). A static right-hand type is its own sole concretisation, so sharing
is trivial and two consistent subtyping checks remain (\ruleName{usub-static}). For a union on
the right, inversion sends each left witness into one branch: the same branch
(\ruleName{usub-union-left}, \ruleName{usub-union-right}), or different branches, in which case
the witnesses no longer interact (\ruleName{usub-union-split}). A constructor on the right must
head both left types, or the case refutes. Tuples proceed componentwise (\ruleName{usub-tuple});
lists and dictionaries need one element witness equivalent to both element types
(\ruleName{usub-list}, \ruleName{usub-dict}), and callable arguments one witness below both
argument types (\ruleName{usub-callable}), the liftings $\UEquiv$ and $\USup$ of
Definition~\ref{def:uequiv-usup}. Each appeal to the hypothesis decreases the order.
\end{proof}

\begin{figure}
\flushleft \shadebox{$\Sigma \vdash \USub(\gty{\tau}, \gty{\tau}'; \gty{\sigma})$}
\begin{mathpar}
\small
\inferrule*[lab={\ruleName{usub-any-right}}]
{
  \strut
}
{
  \Sigma \vdash \USub(\gty{\tau}, \gty{\tau}'; \tyAny)
}
\and
\inferrule*[lab={\ruleName{usub-any-left}}]
{
  \Sigma \vdash \gty{\tau} \gsubty \gty{\sigma}
}
{
  \Sigma \vdash \USub(\tyAny, \gty{\tau}; \gty{\sigma})
}
\and
\inferrule*[lab={\ruleName{usub-never-left}}]
{
  \Sigma \vdash \gty{\tau} \gsubty \gty{\sigma}
}
{
  \Sigma \vdash \USub(\tyNever, \gty{\tau}; \gty{\sigma})
}
\and
\inferrule*[lab={\ruleName{usub-union}}]
{
  \Sigma \vdash \USub(\gty{\tau}, \tyUnion{\gty{\tau}'}{\gty{\tau}''}; \gty{\sigma})
}
{
  \Sigma \vdash \USub(\tyUnion{\gty{\tau}}{\gty{\tau}'}, \gty{\tau}''; \gty{\sigma})
}
\and
\inferrule*[lab={\ruleName{usub-static}}]
{
  \Sigma \vdash \gty{\tau} \gsubty \sigma \\
  \Sigma \vdash \gty{\tau}' \gsubty \sigma
}
{
  \Sigma \vdash \USub(\gty{\tau}, \gty{\tau}'; \sigma)
}
\and
\inferrule*[lab={\ruleName{usub-union-left}}]
{
  \Sigma \vdash \USub(\gty{\tau}, \gty{\tau}'; \gty{\sigma})
}
{
  \Sigma \vdash \USub(\gty{\tau}, \gty{\tau}'; \tyUnion{\gty{\sigma}}{\gty{\sigma}'})
}
\and
\inferrule*[lab={\ruleName{usub-union-right}}]
{
  \Sigma \vdash \USub(\gty{\tau}, \gty{\tau}'; \gty{\sigma}')
}
{
  \Sigma \vdash \USub(\gty{\tau}, \gty{\tau}'; \tyUnion{\gty{\sigma}}{\gty{\sigma}'})
}
\and
\inferrule*[lab={\ruleName{usub-union-split}}]
{
  \Sigma \vdash \gty{\tau} \gsubty \gty{\sigma} \\
  \Sigma \vdash \gty{\tau}' \gsubty \gty{\sigma}'
}
{
  \Sigma \vdash \USub(\gty{\tau}, \gty{\tau}'; \tyUnion{\gty{\sigma}}{\gty{\sigma}'})
}
\and
\inferrule*[lab={\ruleName{usub-list}}]
{
  \Sigma \vdash \UEquiv(\gty{\tau}, \gty{\tau}'; \gty{\sigma})
}
{
  \Sigma \vdash \USub(\tyList{\gty{\tau}}, \tyList{\gty{\tau}'}; \tyList{\gty{\sigma}})
}
\and
\inferrule*[lab={\ruleName{usub-dict}}]
{
  \Sigma \vdash \UEquiv(\gty{\tau}, \gty{\tau}'; \gty{\sigma})
}
{
  \Sigma \vdash \USub(\tyDict{\gty{\tau}}, \tyDict{\gty{\tau}'}; \tyDict{\gty{\sigma}})
}
\and
\inferrule*[lab={\ruleName{usub-tuple}}]
{
  \Sigma \vdash \USub(\gty{\tau}_i, \gty{\tau}'_i; \gty{\sigma}_i) \quad (\forall i)
}
{
  \Sigma \vdash \USub(\tyTuple{\gty{\tau}_1, \ldots, \gty{\tau}_n}, \tyTuple{\gty{\tau}'_1, \ldots, \gty{\tau}'_n}; \tyTuple{\gty{\sigma}_1, \ldots, \gty{\sigma}_n})
}
\and
\inferrule*[lab={\ruleName{usub-callable}}]
{
  \Sigma \vdash \USup(\gty{\tau}_i, \gty{\tau}'_i; \gty{\sigma}_i) \quad (\forall i) \\
  \Sigma \vdash \USub(\gty{\tau}, \gty{\tau}'; \gty{\sigma})
}
{
  \Sigma \vdash \USub(\tyCallable{\gty{\tau}_1, \ldots, \gty{\tau}_n}{\gty{\tau}}, \tyCallable{\gty{\tau}'_1, \ldots, \gty{\tau}'_n}{\gty{\tau}'}; \tyCallable{\gty{\sigma}_1, \ldots, \gty{\sigma}_n}{\gty{\sigma}})
}
\end{mathpar}
\caption{$\USub$: rules extracted from the decidability proof, up to symmetry in the left-hand types}
\label{fig:graduality-usub}
\end{figure}


\begin{definition}[$\UEquiv$ and $\USup$]
\label{def:uequiv-usup}
$\Sigma \vdash \UEquiv(\gty{\tau}, \gty{\tau}'; \gty{\sigma})$ if and only if
$\Sigma \vdash \tau \tyEquiv \sigma$ and $\Sigma \vdash \tau' \tyEquiv \sigma$ for some
$\tau \in \gamma(\gty{\tau})$, $\tau' \in \gamma(\gty{\tau}')$ and $\sigma \in \gamma(\gty{\sigma})$.
$\Sigma \vdash \USup(\gty{\tau}, \gty{\tau}'; \gty{\sigma})$ if and only if
$\Sigma \vdash \sigma \subty \tau$ and $\Sigma \vdash \sigma \subty \tau'$ for some such $\tau$,
$\tau'$ and $\sigma$.
\end{definition}

\noindent The characterisation of $\UEquiv$ by the same method is deferred.

\subsection{Rules for $\USup$}

The rules of \figref{graduality-usub} appeal to $\USup$ at callable arguments, so the method must
be repeated for $\USup$. The case analysis closes every case but one; the positive cases give
\figref{graduality-usup}.

The induction is on the total size of the triple; cases as before. On the right, $\tyAny$ admits
the witness $\tyNever$, below every type (\ruleName{usup-any-right}); a static right-hand type is
its own sole concretisation, so sharing is trivial and two consistent subtyping checks remain
(\ruleName{usup-static}); a right-hand union that is not static is the residual case, taken up
below. In the remaining cases the shared witness is constructor-headed. A left type $\tyAny$ or
$\tyObject$ takes the witness $\tyObject$, above every type, leaving a consistent subtyping check
on the other left type (\ruleName{usup-any-left}, \ruleName{usup-object-left}). For a union on the
left, inversion sends the shared witness below one component, and the witness for the other
component is unconstrained (\ruleName{usup-union-left}, \ruleName{usup-union-right}). A
constructor on the right must head both left types, or the case refutes. Tuples proceed
componentwise (\ruleName{usup-tuple}); lists and dictionaries need one element witness equivalent
to both element types (\ruleName{usup-list}, \ruleName{usup-dict}); callable arguments need one
argument witness above both argument types, instances of $\USub$ (\ruleName{usup-callable}), so
the two predicates are mutually recursive.

\begin{figure}
\flushleft \shadebox{$\Sigma \vdash \USup(\gty{\tau}, \gty{\tau}'; \gty{\sigma})$}
\begin{mathpar}
\small
\inferrule*[lab={\ruleName{usup-any-right}}]
{
  \strut
}
{
  \Sigma \vdash \USup(\gty{\tau}, \gty{\tau}'; \tyAny)
}
\and
\inferrule*[lab={\ruleName{usup-any-left}}]
{
  \Sigma \vdash \gty{\sigma} \gsubty \gty{\tau}
}
{
  \Sigma \vdash \USup(\tyAny, \gty{\tau}; \gty{\sigma})
}
\and
\inferrule*[lab={\ruleName{usup-object-left}}]
{
  \Sigma \vdash \gty{\sigma} \gsubty \gty{\tau}
}
{
  \Sigma \vdash \USup(\tyObject, \gty{\tau}; \gty{\sigma})
}
\and
\inferrule*[lab={\ruleName{usup-static}}]
{
  \Sigma \vdash \sigma \gsubty \gty{\tau} \\
  \Sigma \vdash \sigma \gsubty \gty{\tau}'
}
{
  \Sigma \vdash \USup(\gty{\tau}, \gty{\tau}'; \sigma)
}
\and
\inferrule*[lab={\ruleName{usup-union-left}}]
{
  \Sigma \vdash \USup(\gty{\tau}, \gty{\tau}''; \gty{\sigma})
}
{
  \Sigma \vdash \USup(\tyUnion{\gty{\tau}}{\gty{\tau}'}, \gty{\tau}''; \gty{\sigma})
}
\and
\inferrule*[lab={\ruleName{usup-union-right}}]
{
  \Sigma \vdash \USup(\gty{\tau}', \gty{\tau}''; \gty{\sigma})
}
{
  \Sigma \vdash \USup(\tyUnion{\gty{\tau}}{\gty{\tau}'}, \gty{\tau}''; \gty{\sigma})
}
\and
\inferrule*[lab={\ruleName{usup-list}}]
{
  \Sigma \vdash \UEquiv(\gty{\tau}, \gty{\tau}'; \gty{\sigma})
}
{
  \Sigma \vdash \USup(\tyList{\gty{\tau}}, \tyList{\gty{\tau}'}; \tyList{\gty{\sigma}})
}
\and
\inferrule*[lab={\ruleName{usup-dict}}]
{
  \Sigma \vdash \UEquiv(\gty{\tau}, \gty{\tau}'; \gty{\sigma})
}
{
  \Sigma \vdash \USup(\tyDict{\gty{\tau}}, \tyDict{\gty{\tau}'}; \tyDict{\gty{\sigma}})
}
\and
\inferrule*[lab={\ruleName{usup-tuple}}]
{
  \Sigma \vdash \USup(\gty{\tau}_i, \gty{\tau}'_i; \gty{\sigma}_i) \quad (\forall i)
}
{
  \Sigma \vdash \USup(\tyTuple{\gty{\tau}_1, \ldots, \gty{\tau}_n}, \tyTuple{\gty{\tau}'_1, \ldots, \gty{\tau}'_n}; \tyTuple{\gty{\sigma}_1, \ldots, \gty{\sigma}_n})
}
\and
\inferrule*[lab={\ruleName{usup-callable}}]
{
  \Sigma \vdash \USub(\gty{\tau}_i, \gty{\tau}'_i; \gty{\sigma}_i) \quad (\forall i) \\
  \Sigma \vdash \USup(\gty{\tau}, \gty{\tau}'; \gty{\sigma})
}
{
  \Sigma \vdash \USup(\tyCallable{\gty{\tau}_1, \ldots, \gty{\tau}_n}{\gty{\tau}}, \tyCallable{\gty{\tau}'_1, \ldots, \gty{\tau}'_n}{\gty{\tau}'}; \tyCallable{\gty{\sigma}_1, \ldots, \gty{\sigma}_n}{\gty{\sigma}})
}
\end{mathpar}
\caption{$\USup$: rules extracted from the case analysis, up to symmetry in the left-hand types; a right-hand union that is not static has no rule}
\label{fig:graduality-usup}
\end{figure}


The residual case is a right-hand union $\tyUnion{\gty{\sigma}}{\gty{\sigma}'}$ that is not
static. Its witness is a union, below a type exactly when both components are, so the two
constraints become four, relating each component witness to each left witness; every witness now
appears in two constraints. The case amounts to $\USub(\gty{\sigma}, \gty{\sigma}'; \gty{\tau})$
and $\USub(\gty{\sigma}, \gty{\sigma}'; \gty{\tau}')$ with the component witnesses shared between
the two instances, which no predicate over three gradual types expresses. Splitting the union
instead, with one $\USup$ premise per component, is unsound:
$\USup(\tyList{\tyAny}, \tyList{\tyAny}; \tyList{\tyInt})$ and
$\USup(\tyList{\tyAny}, \tyList{\tyAny}; \tyList{\tyStr})$ hold, but
$\USup(\tyList{\tyAny}, \tyList{\tyAny}; \tyUnion{\tyList{\tyInt}}{\tyList{\tyStr}})$ needs an
element witness equivalent to both $\tyInt$ and $\tyStr$.

Sharing the pair of component witnesses suggests a predicate over four gradual types, but the
widening does not stop there: a union on the left then sends each component witness below a
component of the left union independently, so the constraints no longer relate every pair. The
general form is a finite set of subtyping and equivalence constraints over shared witnesses, a
constraint system in the sense of \citet{garcia2015}; its treatment is deferred.
